\documentclass[11pt]{article}
\usepackage[margin=1in]{geometry}
\usepackage{amsmath,amssymb,amsthm,mathtools}
\usepackage{booktabs}
\usepackage{enumitem}
\usepackage{hyperref}
\usepackage{xcolor}
\hypersetup{colorlinks=true,linkcolor=blue,citecolor=blue,urlcolor=blue}

\newtheorem{theorem}{Theorem}[section]
\newtheorem{lemma}[theorem]{Lemma}
\newtheorem{proposition}[theorem]{Proposition}
\newtheorem{corollary}[theorem]{Corollary}
\newtheorem{definition}[theorem]{Definition}
\newtheorem{example}[theorem]{Example}
\newtheorem{remark}[theorem]{Remark}

\DeclareMathOperator{\Range}{Ran}
\DeclareMathOperator{\Null}{Ker}
\DeclareMathOperator{\tr}{tr}
\DeclareMathOperator{\diag}{diag}
\newcommand{\XiC}{\Xi_C}
\newcommand{\KLL}{K_{LL}}
\newcommand{\KDL}{K_{DL}}
\newcommand{\KLD}{K_{LD}}
\newcommand{\KDD}{K_{DD}}
\newcommand{\KL}{K_L}
\newcommand{\KD}{K_D}
\newcommand{\Elegal}{E_C}

\title{Adequacy Residuals and Blind-Spot Currency\\
A Standalone \texorpdfstring{$\Xi$}{Xi}-Theory for Formed-Layer Membranes}
\author{Six Birds / RATCHET working draft}
\date{Step 57 draft scaffold}

\begin{document}
\maketitle

\begin{abstract}
This note isolates the adequacy residual
\[
\Xi_C(D\mid L)=K_{DD}-K_{DL}K_{LL}^{\dagger}K_{LD}
\]
as a standalone object in formed-layer membrane theory.  The native probe family
\(L\) records what a formed layer officially sees, while the dissolving probe family
\(D\) records ways the layer could be broken or reduced to hidden needles.  The
matrix \(\Xi_C(D\mid L)\) is the optimal residual currency: it is exactly the part
of the dissolving readout not explained by native probes in the legal energy-dual
geometry.  The note proves the projection formula, optimal residual formula,
exact-adequacy criterion, strict-extension monotonicity, chain rule, predictive
residual ladder, and membrane-promotion theorem.  It also records the necessary
null-mode and nonclaim boundaries.  The scope is closure-only: the theory applies
to formed carriers and declared probe families, not to arbitrary raw substrates.
\end{abstract}

\section{Why \texorpdfstring{$\Xi$}{Xi} deserves a standalone theory}
The larger formed-layer membrane theory says that a closure/layer is protected
when every native predictive probe recombination has a finite audited price.  That
is a native statement.  It does not automatically imply that every layer-dissolving
probe is priced.  The missing bridge is adequacy.

The adequacy residual
\[
\Xi_C(D\mid L)
\]
answers the question:
\begin{quote}
What can a dissolving probe family \(D\) see that the native family \(L\) cannot
see, after legal energy quotienting?
\end{quote}
It is therefore the layer's blind-spot currency.  If \(\Xi=0\), the native probes
are adequate for \(D\).  If \(\Xi\succ0\), there is a legal hidden direction or a
positive residual budget that must be paid before a native membrane can be
promoted to a dissolving membrane.

\section{Legal quotient and currency matrices}
Let \(E\) be a finite-dimensional Hilbert space and let \(C=C^*\succeq0\) be an
audit form.  Define the legal energy quotient
\[
E_C=(\ker C)^\perp,
\qquad
C_0=C|_{E_C}>0.
\]
Let \(L:E\to Y\) be the native probe family and \(D:E\to Z\) the dissolving probe
family.  We say \(L,D\) are null-legal if
\[
L|_{\ker C}=0,
\qquad
D|_{\ker C}=0.
\]
Write \(L_0=L|_{E_C}\) and \(D_0=D|_{E_C}\).

\begin{definition}[Currency blocks]
For null-legal \(L,D\), define
\[
K_{LL}=L_0C_0^{-1}L_0^*,
\qquad
K_{DL}=D_0C_0^{-1}L_0^*,
\]
\[
K_{LD}=K_{DL}^*,
\qquad
K_{DD}=D_0C_0^{-1}D_0^*.
\]
\end{definition}

\begin{definition}[Adequacy residual / blind-spot currency]
The adequacy residual of \(D\) relative to \(L\) is
\[
\boxed{
\Xi_C(D\mid L)=K_{DD}-K_{DL}K_{LL}^{\dagger}K_{LD}.
}
\]
\end{definition}

\begin{remark}[Null-mode warning]
If \(D\) sees a vector in \(\ker C\), then the true variational capacity of a
corresponding dissolving probe is infinite.  A pseudoinverse expression that
silently assigns zero to that null component is not an accepted currency record.
\end{remark}

\section{Projection and Schur geometry}
Define
\[
T_L=L_0C_0^{-1/2}:E_C\to Y,
\qquad
T_D=D_0C_0^{-1/2}:E_C\to Z.
\]
Let \(P_L\) be the orthogonal projection of \(E_C\) onto \(\overline{\Range(T_L^*)}\)
(which is simply \(\Range(T_L^*)\) in finite dimension).

\begin{theorem}[Projection formula]
\[
\boxed{
\Xi_C(D\mid L)=T_D(I-P_L)T_D^*\succeq0.
}
\]
\end{theorem}

\begin{proof}
The block matrix
\[
\begin{pmatrix}
K_{LL}&K_{LD}\\
K_{DL}&K_{DD}
\end{pmatrix}
=
\begin{pmatrix}
T_L\\ T_D
\end{pmatrix}
\begin{pmatrix}
T_L\\ T_D
\end{pmatrix}^*
\]
is a Gram matrix.  The Schur complement of \(K_{LL}=T_LT_L^*\) is the squared
norm of the component of \(T_D^*\) orthogonal to \(\Range(T_L^*)\).  Algebraically,
\(T_L^*(T_LT_L^*)^\dagger T_L=P_L\), hence
\[
K_{DD}-K_{DL}K_{LL}^\dagger K_{LD}
=T_D(I-P_L)T_D^*.
\]
\end{proof}

\begin{corollary}[Exact adequacy criterion]
The following are equivalent:
\[
\Xi_C(D\mid L)=0,
\]
\[
D_0=A L_0\quad\text{for some }A:Y\to Z,
\]
\[
D_0=A_*L_0\quad\text{where }A_*=K_{DL}K_{LL}^{\dagger},
\]
\[
\ker L_0\subseteq\ker D_0.
\]
\end{corollary}

\begin{proof}
\(\Xi=0\) iff \(T_D^*\) lands in \(\Range(T_L^*)\), equivalently every energy-dual
linear functional generated by \(D\) is generated by \(L\).  This is equivalent to
factorization through \(L_0\), which is equivalent to \(\ker L_0\subseteq\ker D_0\)
in finite dimension.  The minimizing factor is the least-squares factor
\(A_*=K_{DL}K_{LL}^{\dagger}\).
\end{proof}

\section{Optimal residual and minimum blind-spot theorem}
\begin{theorem}[Optimal residual identity]
For any \(A:Y\to Z\),
\[
\boxed{
(D_0-AL_0)C_0^{-1}(D_0-AL_0)^*
=
\Xi_C(D\mid L)+(A-A_*)K_{LL}(A-A_*)^*.
}
\]
Consequently \(\Xi_C(D\mid L)\) is the Loewner-minimal residual currency among
all factorizations \(D\approx AL\).
\end{theorem}

\begin{proof}
Expand the left-hand side using the currency blocks:
\[
K_{DD}-AK_{LD}-K_{DL}A^*+AK_{LL}A^*.
\]
Completing the square around \(A_*=K_{DL}K_{LL}^{\dagger}\), and using the range
relations valid for Gram matrices, gives the stated identity.
\end{proof}

\begin{remark}[Interpretation]
\(\Xi\) is not a modeling residual chosen by hand.  It is the unique optimal
blind-spot currency.  It is what remains of the dissolving family after the best
native reconstruction is subtracted.
\end{remark}

\section{Promotion from native membrane to dissolving membrane}
\begin{theorem}[Adequacy promotion]
If
\[
K_{LL}\preceq \Theta_Y
\qquad\text{and}\qquad
\Xi_C(D\mid L)\preceq \Xi_Z,
\]
then
\[
\boxed{
K_{DD}\preceq A_*\Theta_YA_*^*+\Xi_Z.
}
\]
In particular, if \(\Xi_C(D\mid L)=0\), then native membrane control transfers
exactly to the dissolving family through \(A_*\).
\end{theorem}

\begin{proof}
By Schur decomposition,
\[
K_{DD}=K_{DL}K_{LL}^{\dagger}K_{LD}+\Xi
=A_*K_{LL}A_*^*+\Xi.
\]
Use \(K_{LL}\preceq\Theta_Y\) and \(\Xi\preceq\Xi_Z\).
\end{proof}

\section{Strict extension and chain rule}
Let \(M:E\to W\) be a new native probe family.  Write
\[
L^+=\begin{bmatrix}L\\M\end{bmatrix}.
\]
Define conditional currency blocks after conditioning on \(L\):
\[
K_{MM\mid L}=K_{MM}-K_{ML}K_{LL}^{\dagger}K_{LM},
\]
\[
K_{DM\mid L}=K_{DM}-K_{DL}K_{LL}^{\dagger}K_{LM}.
\]

\begin{theorem}[Strict-extension monotonicity and chain rule]
\[
\boxed{
\Xi_C(D\mid L^+)
=
\Xi_C(D\mid L)-K_{DM\mid L}K_{MM\mid L}^{\dagger}K_{MD\mid L}.
}
\]
Hence
\[
\boxed{
\Xi_C(D\mid L^+)\preceq \Xi_C(D\mid L).
}
\]
Strict decrease occurs exactly in the directions of the old residual seen by the
new native family \(M\).
\end{theorem}

\begin{proof}
Apply the Schur-complement associativity identity to the Gram matrix of
\((L,M,D)C_0^{-1/2}\).  Conditioning first on \(L\), then on \(M\) within the
orthogonal complement of \(L\), yields the stated decrement.
\end{proof}

\begin{corollary}[Same-family saturation]
If \(M=B L\) for some \(B:Y\to W\), then
\[
K_{MM\mid L}=0,
\qquad
\Xi_C(D\mid [L;M])=\Xi_C(D\mid L).
\]
Repeating, relabeling, or post-processing the old native family does not reduce
blind-spot currency.
\end{corollary}

\section{Predictive adequacy residual ladders}
At stage \(j\), let
\[
(E_j,C_j,L_j,D_j,R^Y_j,R^Z_j)
\]
be a predictive membrane record.  Transported residuals live in a common
dissolving response space \(Z_\ast\):
\[
\widehat \Xi_j=R^Z_j\Xi_{C_j}(D_j\mid L_j)(R^Z_j)^*.
\]

\begin{theorem}[Summable residual propagation]
Suppose
\[
\widehat\Xi_{j+1}\preceq (1+\varepsilon_j)B_j\widehat\Xi_jB_j^*+F_j,
\]
and in a common response space \(B_j=I\),
\[
\prod_j(1+\varepsilon_j)<\infty,
\qquad
\sum_jF_j\preceq F_\infty.
\]
Then all stages satisfy
\[
\widehat\Xi_n\preceq P_\infty(\widehat\Xi_0+F_\infty)
\]
for a finite constant \(P_\infty\).  Thus a predictive adequacy budget persists
under summable strict-extension defects.
\end{theorem}

\begin{remark}[Non-summable failure]
Finite adequacy residual at every finite stage does not imply predictive
adequacy.  Residual increments of order \(1/j\) can diverge in total.
\end{remark}

\section{Examples and application templates}
\subsection{Blind-spot counterexample}
Let
\[
E=\mathbb R^2,
\quad C=I,
\quad L(x_1,x_2)=x_1,
\quad D(x_1,x_2)=Mx_2.
\]
Then \(K_{LL}=1\), but
\[
\Xi_C(D\mid L)=M^2.
\]
A native membrane can be perfectly bounded while the dissolving family has an
arbitrarily large blind spot.

\subsection{Navier adequacy template}
For a Fourier cutoff carrier with Sobolev audit \(C_s\), native family \(L\) may
record energy/enstrophy or shell features, while \(D\) records pointwise velocity,
vorticity, or strain concentration.  The residual \(\Xi_{C_s}(D\mid L)\) measures
the blind spot between the energy ledger and blow-up probes.  In dimension three,
classical energy/enstrophy does not uniformly price pointwise gradient probes; a
strict extension must add higher coercivity, channelization, parabolic smoothing,
source-admission constraints, or explicit nonclaims.

\subsection{Security and immune analogies}
In cryptography, \(L\) is a declared adversary model and \(D\) an expanded attack
family.  Then \(\Xi\) is the unknown-attack blind spot.  In immunology, \(L\) is
known antigen recognition and \(D\) novel pathogen presentation.  The same
Schur-residual mathematics applies; only the carrier and audit change.

\section{Nonclaims and scope}
This theory does not claim that every raw substrate is needle-free.  It claims
that a formed closure with declared native and dissolving probe families has a
measurable adequacy residual.  A zero or controlled residual licenses promotion
from native membrane to dissolving membrane only within the declared audit,
transport, and null-mode legality records.

\begin{center}
\fbox{\parbox{0.91\linewidth}{
\textbf{Slogan.}  \(\Xi_C(D\mid L)\) is what the layer does not know about its own possible dissolution.  Strict extension reduces this blind spot exactly when the new native probes see the old residual.
}}
\end{center}

\end{document}
